\begin{statement}{C5a}{Proved}{The normalization quotient and its cohomology map}
Let $Y$ be a smooth compact complex threefold, $S\subset Y$ a reduced normal-crossing divisor, and $\nu:\widetilde S\to Y$ the normalization followed by inclusion. This is a finite holomorphic immersion on each local branch. Put $\mathcal T=T_Y(-\log S)$. There is a canonical exact sequence of coherent $\OO_Y$-modules
\[
0\longrightarrow\mathcal T\longrightarrow T_Y
 \longrightarrow\nu_*N_\nu\longrightarrow0,
\qquad N_\nu=K_{\widetilde S}\otimes\nu^*K_Y^{-1}.      \tag{C.3}
\]
For every holomorphic line bundle $L$ on $Y$, the image of the natural cohomology map is exactly
\[
\begin{split}
\im\bigl(H^2(Y,T_Y\otimes L)&\longrightarrow
 H^2(\widetilde S,N_\nu\otimes\nu^*L)\bigr)\\
 &=\ker\bigl(H^2(\widetilde S,N_\nu\otimes\nu^*L)
 \longrightarrow H^3(Y,\mathcal T\otimes L)\bigr).
\end{split}                                             \tag{C.4}
\]
The target identifies by Serre duality with
$H^0(\widetilde S,\nu^*(K_Y\otimes L^{-1}))^\vee$.
\end{statement}
\begin{proof}
In coordinates with $S=(z_1\cdots z_r=0)$, $1\le r\le3$, the map sends a vector field $\sum a_i\partial_{z_i}$ to the normal coefficient $a_i|_{z_i=0}$ on each branch $i\le r$. The separate coefficients lift independently to holomorphic functions, making the map surjective. Its kernel consists precisely of fields with $z_i\mid a_i$ for $i\le r$, which is $T_Y(-\log S)$. Coordinate changes act on each coefficient through the normal-line transition and permute the branches; hence the map and its exactness descend canonically to the normalization. The immersion tangent sequence on $\widetilde S$ gives $\det\nu^*T_Y=\det T_{\widetilde S}\otimes N_\nu$, proving the adjunction expression. Tensoring by the locally free line bundle $L$ preserves exactness. Finite pushforward has no higher direct images, so the long exact sequence in coherent cohomology gives (C.4). The smooth compact surface $\widetilde S$ has the usual Serre pairing, which gives the stated dual target. In particular the vanishing of $H^3(Y,\mathcal T\otimes L)$ makes this cohomology map surjective; no such vanishing has been suppressed in (C.4).
\end{proof}


% C5b from part_C5.tex

In this subsection $X=X_c$ is the period construction of \cite{AC26}, compared with the quotient construction of \cite{Engel2026}, with admissible parameter $c$, standard clockwise affine vectors
$v_1=(1,2,-4,0)^t$, $v_2=(-1,-3,3,0)^t$, and untwisted cusp collar. Thus $p=12l_0-4l_1-3l_2=-1$. We write $\mathfrak B=\PP(3,4)$, $\rho:\mathfrak B\to B$ for its coarse map, and $\bar f:X\to\mathfrak B$ for the root factorization. Let $\mathcal T=T_X(-\log W)$ and let $\nu:\mathrm dP_6\to X$ normalize $W$. The root generator is fixed by
\[
 M_n=\bar f^*\OO_{\mathfrak B}(n)
     =\OO_X(n(S_1-S_2)),\quad n\in\Z.
\]
The exponential isomorphism $H^1(X,\OO_X)\to\Pic(X)$ identifies $\Pic(X)$ with $\C$. We scale its coordinate so that $M_z=\exp(z\gamma^{0,1})$ and $M_1=\OO_X(S_1-S_2)$. The sign of this equality is fixed by the clockwise transgression $12\gamma\mapsto-p=1$. This is checked by evaluating the two finite lifts on $v_1,v_2$: their integral periods are $4l_1,3l_2$, while the cusp contributes $-12l_0$, whose sum is $-p$. The exponential sequence therefore gives $\exp(12\gamma^{0,1})=f^*\OO_B(1)$. Its injectivity and the root relation give the stated $M_1$ equality.

